\frfilename{mt253.tex}
\versiondate{18.4.08}

\def\chaptername{Ukuran produk}
\def\sectionname{Produk tensor}

\newsection{253}

Teorema-teorema pada bagian terakhir menunjukkan bahwa fungsi terintegralkan pada
produk dua ruang ukur dapat dikaji secara efektif melalui
integrasi pada setiap ruang faktor secara terpisah. Dalam bagian ini saya menyajikan
hubungan yang sangat mencolok antara ruang $L^1$ dari suatu ukuran
produk dan ruang-ruang $L^1$ dari faktor-faktornya, yang bahkan menentukan
$L^1$ produk hingga isomorfisme sebagai kisi Banach. Saya memulai dengan
catatan singkat tentang operator bilinear (253A) dan uraian tentang operator
bilinear kanonik dari $L^1(\mu)\times L^1(\nu)$ ke $L^1(\mu\times\nu)$
(253B-253E). Teorema utama bagian ini adalah 253F, yang menunjukkan bahwa
peta kanonik ini bersifat universal bagi operator bilinear kontinu dari
$L^1(\mu)\times L^1(\nu)$ ke ruang Banach; peta itu juga menentukan
urutan pada $L^1(\mu\times\nu)$ (253G). Saya mengakhiri bagian ini dengan uraian tentang
suatu jenis mendasar operator ekspektasi bersyarat (253H) dan catatan
tentang produk ukuran integral tak tentu (253I) serta integral atas
untuk jenis-jenis fungsi khusus (253J, 253K).

\leader{253A}{Operator bilinear}\cmmnt{ Sebelum menelaah bagian mana pun dari
teori ukuran dalam bagian ini, saya memperkenalkan suatu konsep dari teori
ruang linear.

\medskip

} {\bf (a)} Misalkan $U$, $V$, dan $W$ adalah ruang linear atas
$\Bbb  R$\cmmnt{ (atau, sesungguhnya, atas sebarang lapangan)}. Suatu peta
$\phi:U\times V\to W$ disebut {\bf bilinear} jika linear dalam setiap variabel
secara terpisah, yaitu,

\Centerline{$\phi(u_1+u_2,v)=\phi(u_1,v)+\phi(u_2,v)$,}

\Centerline{$\phi(u,v_1+v_2)=\phi(u,v_1)+\phi(u,v_2)$,}

\Centerline{$\phi(\alpha u,v)=\alpha\phi(u,v)=\phi(u,\alpha v)$}

\noindent untuk setiap $u$, $u_1$, $u_2\in U$, $v$, $v_1$, $v_2\in V$, dan
skalar $\alpha$. Perhatikan bahwa $\phi$ semacam itu menghasilkan, dan pada
gilirannya dapat didefinisikan oleh, suatu operator linear
$T:U\to\eurm L(V;W)$\cmmnt{, dengan menuliskan
$\eurm L(V;W)$ untuk ruang operator linear dari $V$ ke $W$}, dengan

\Centerline{$(Tu)(v)=\phi(u,v)$}

\noindent untuk setiap $u\in U$, $v\in V$. \cmmnt{Oleh karena itu, atau dengan cara lain, kita
dapat melihat, misalnya, bahwa} $\phi(0,v)=\phi(u,0)=0$ kapan pun $u\in U$ dan
$v\in V$.

Jika $W'$ adalah ruang linear lain\cmmnt{ atas lapangan yang sama,} dan
$S:W\to W'$ suatu operator linear, maka $S\phi:U\times V\to W'$ bersifat
bilinear.

\spheader 253Ab Sekarang andaikan $U$, $V$, dan $W$ adalah
ruang bernorma, dan $\phi:U\times V\to W$ suatu operator bilinear.
\cmmnt{Maka kita
mengatakan bahwa} $\phi$ {\bf terbatas} jika
$\sup\{\|\phi(u,v)\|:\|u\|\le 1,\,\|v\|\le 1\}$ berhingga, dan dalam hal ini
kita menyebut supremum tersebut sebagai
norma $\|\phi\|$ dari $\phi$. \cmmnt{Perhatikan bahwa}
$\|\phi(u,v)\|\le\|\phi\|\|u\|\|v\|$ untuk setiap $u\in U$, $v\in
V$\prooflet{ (karena

\Centerline{$\|\phi(u,v)\|=\alpha\beta\|\phi(\alpha^{-1}u,\beta^{-1}v)\|
\le\alpha\beta\|\phi\|$}

\noindent kapan pun $\alpha>\|u\|$, $\beta>\|v\|$)}.

Jika $W'$ adalah ruang bernorma lain dan $S:W\to W'$ suatu operator linear
terbatas, maka $S\phi:U\times V\to W'$ adalah operator bilinear terbatas,
dan $\|S\phi\|\le\|S\|\|\phi\|$.

\leader{253B}{Definisi}\cmmnt{ Operator bilinear terpenting dalam
bagian ini didasarkan pada gagasan berikut.} Misalkan $f$ dan $g$ adalah
fungsi bernilai real. Saya akan menuliskan $f\otimes g$ untuk fungsi
$(x,y)\mapsto f(x)g(y):\dom f\times\dom g\to\Bbb R$.

\leader{253C}{Proposisi}
(a) Misalkan $X$ dan $Y$ adalah himpunan, dan $\Sigma$, $\Tau$ masing-masing
aljabar-sigma dari subhimpunan $X$, $Y$. Jika $f$ adalah fungsi bernilai real
yang terukur terhadap $\Sigma$ dan didefinisikan pada suatu subhimpunan $X$,
serta $g$ adalah fungsi bernilai real yang terukur terhadap $\Tau$ dan
didefinisikan pada suatu subhimpunan $Y$, maka $f\otimes g$\cmmnt{, sebagaimana
didefinisikan dalam 253B,} terukur terhadap $\Sigma\tensorhat\Tau$.

(b) Misalkan $(X,\Sigma,\mu)$ dan $(Y,\Tau,\nu)$ adalah ruang ukur, dan
$\lambda$ ukuran produk c.l.d.\ pada $X\times Y$. Jika
$f\in\eusm L^0(\mu)$ dan $g\in\eusm L^0(\nu)$, maka
$f\otimes g\in\eusm L^0(\lambda)$.

\cmmnt{\medskip

\noindent{\bf Catatan} Ingat dari 241A bahwa $\eusm L^0(\mu)$ adalah
ruang fungsi bernilai real yang terukur secara virtual terhadap $\mu$ dan
didefinisikan pada subhimpunan koterabaikan terhadap $\mu$ dari $X$.}

\proof{{\bf (a)} Intinya ialah bahwa $f\otimes\chi Y$ terukur terhadap
$\Sigma\tensorhat\Tau$, sebab untuk setiap $\alpha\in\Bbb R$
terdapat $E\in\Sigma$ sedemikian sehingga

\Centerline{$\{x:f(x)\ge\alpha\}=E\cap\dom f$,}

\noindent sehingga

\Centerline{$\{(x,y):(f\otimes\chi Y)(x,y)\ge\alpha\}
=(E\cap\dom f)\times Y=(E\times Y)\cap\dom(f\otimes\chi Y)$,}

\noindent dan tentu saja $E\times Y\in\Sigma\tensorhat\Tau$. Dengan cara yang sama,
$\chi X\otimes g$ terukur terhadap $\Sigma\tensorhat\Tau$ dan
$f\otimes g=(f\otimes\chi Y)\times(\chi X\otimes g)$ terukur terhadap
$\Sigma\tensorhat\Tau$.

\medskip

{\bf (b)} Ambil subhimpunan koterabaikan $E\in\Sigma$, $F\in\Tau$ dari $X$,
$Y$ berturut-turut, sedemikian sehingga $E\subseteq\dom f$, $F\subseteq\dom g$,
$f\restr E$ terukur terhadap $\Sigma$ dan $g\restr F$ terukur terhadap $\Tau$.
Tuliskan $\Lambda$ untuk domain $\lambda$. Maka
$\Sigma\tensorhat\Tau\subseteq\Lambda$ (251Ia). Selain itu, $E\times F$
koterabaikan terhadap $\lambda$, karena

$$\eqalign{\lambda((X\times Y)\setminus(E\times F))
&\le\lambda((X\setminus E)\times Y)+\lambda(X\times(Y\setminus F))\cr
&=\mu(X\setminus E)\cdot\nu Y+\mu X\cdot\nu(Y\setminus F)
=0\cr}$$

\noindent (juga dari 251Ia). Jadi $\dom(f\otimes g)\supseteq E\times F$
koterabaikan. Selain itu, menurut (a),
$(f\otimes g)\restr(E\times F)
=(f\restr E)\otimes(g\restr F)$ terukur terhadap $\Sigma\tensorhat\Tau$,
dan karenanya terukur terhadap $\Lambda$, serta $f\otimes g$ terukur secara
virtual. Jadi $f\otimes g\in\eusm L^0(\lambda)$, sebagaimana dinyatakan.
}%end of proof of 253C

\leader{253D}{}\cmmnt{ Sekarang kita dapat menerapkan gagasan 253B-253C pada
fungsi terintegralkan.

\medskip

\noindent}{\bf Proposisi} Misalkan $(X,\Sigma,\mu)$ dan
$(Y,\Tau,\nu)$ adalah ruang ukur, dan tuliskan $\lambda$ untuk ukuran
produk c.l.d.\ pada $X\times Y$. Jika $f\in\eusm L^1(\mu)$ dan
$g\in\eusm L^1(\nu)$, maka $f\otimes g\in\eusm L^1(\lambda)$ dan
$\int f\otimes g\,d\lambda=\int f\,d\mu\int g\,d\nu$.

\cmmnt{\medskip

\noindent{\bf Catatan} Saya mengikuti \S242 dengan menuliskan $\eusm L^1(\mu)$ untuk
ruang fungsi bernilai real yang terintegralkan terhadap $\mu$.
}

\proof{{\bf (a)} Tinjau terlebih dahulu
kasus $f=\chi E$, $g=\chi F$ dengan $E\in\Sigma$, $F\in\Tau$ berukuran
berhingga; maka $f\otimes g=\chi(E\times F)$ terintegralkan terhadap $\lambda$ dengan
integral

\Centerline{$\lambda(E\times F)=\mu E\cdot\nu F
=\int f\,d\mu\cdot\int g\,d\nu$,}

\noindent menurut 251Ia.

\medskip

{\bf (b)} Segera diperoleh bahwa $f\otimes g$ adalah fungsi sederhana-$\lambda$,
dengan $\int f\otimes g\,d\lambda=\int f\,d\mu\int g\,d\nu$, kapan pun $f$
merupakan fungsi sederhana-$\mu$ dan $g$ merupakan fungsi sederhana-$\nu$.

\medskip

{\bf (c)} Jika $f$ dan $g$ adalah fungsi terintegralkan nonnegatif, terdapat
barisan tak-menurun $\sequencen{f_n}$, $\sequencen{g_n}$ yang terdiri atas fungsi
sederhana nonnegatif
dan masing-masing konvergen hampir di mana-mana ke $f$, $g$; sekarang
perhatikan bahwa jika $E\subseteq X$, $F\subseteq Y$ koterabaikan, $E\times
F$ koterabaikan dalam $X\times Y$, sebagaimana dicatat dalam pembuktian 253C, sehingga
barisan tak-menurun $\sequencen{f_n\otimes g_n}$ yang terdiri atas fungsi
sederhana-$\lambda$ konvergen hampir di mana-mana ke $f\otimes g$,
dan

\Centerline{$\int f\otimes g\,d\lambda
=\lim_{n\to\infty}\int f_n\otimes g_nd\lambda
=\lim_{n\to\infty}\int f_nd\mu\int g_nd\nu
=\int f\,d\mu\int g\,d\nu$}

\noindent menurut teorema B.Levi.

\medskip

{\bf (d)} Akhirnya, untuk $f$ dan $g$ umum, kita dapat menyatakannya sebagai
selisih $f^+-f^-$, $g^+-g^-$ dari fungsi-fungsi terintegralkan nonnegatif,
dan melihat bahwa

\Centerline{$\int f\otimes g\,d\lambda
=\int f^+\otimes g^+-f^+\otimes g^--f^-\otimes g^+
+f^-\otimes g^-d\lambda
=\int f\,d\mu\int g\,d\nu.$}
}%end of proof of 253D

\leader{253E}{Peta kanonik $L^1\times L^1\to L^1$}\cmmnt{ Saya
melanjutkan argumen dari 253D. Karena $E\times F$ koterabaikan
dalam $X\times Y$ kapan pun $E$ dan $F$ merupakan
subhimpunan koterabaikan dari $X$ dan $Y$,
$f_1\otimes g_1=f\otimes g\,\,\lambda$-hampir di mana-mana kapan pun $f=f_1\,\,\mu$-hampir di mana-mana
dan $g=g_1\,\,\nu$-hampir di mana-mana.} Kita dapat\cmmnt{ karena itu} mendefinisikan
$u\otimes v\in L^1(\lambda)$, untuk $u\in L^1(\mu)$ dan $v\in L^1(\nu)$,
dengan menetapkan bahwa
$u\otimes v=(f\otimes g)^{\ssbullet}$ kapan pun $u=f^{\ssbullet}$ dan
$v=g^{\ssbullet}$. \dvro{P}{

Sekarang jika $f$, $f_1$, $f_2\in\eusm L^1(\mu)$, $g$, $g_1$,
$g_2\in\eusm L^1(\nu)$, dan $a\in\Bbb R$,

\Centerline{$(f_1+f_2)\otimes g=(f_1\otimes g)+(f_2\otimes g)$,}

\Centerline{$f\otimes(g_1+g_2)=(f\otimes g_1)+(f\otimes g_2)$,}

\Centerline{$(af)\otimes g=a(f\otimes g)=f\otimes(ag)$.}

\noindent Segera diperoleh bahwa p}eta $(u,v)\mapsto u\otimes v$ bersifat
bilinear.

\cmmnt{Selain itu, jika $f\in\eusm L^1(\mu)$ dan $g\in\eusm L^1(\nu)$,
$|f|\otimes|g|=|f\otimes g|$, sehingga
$\int|f\otimes g|d\lambda=\int|f|d\mu\int|g|d\nu$. Karena itu}

\Centerline{$\|u\otimes v\|_1=\|u\|_1\|v\|_1$}

\noindent untuk setiap $u\in L^1(\mu)$, $v\in L^1(\nu)$. Secara khusus,
operator bilinear $\otimes$ terbatas, dengan norma $1$ (kecuali dalam kasus sepele
ketika salah satu dari $L^1(\mu)$, $L^1(\nu)$ berdimensi $0$).

\leader{253F}{}\cmmnt{ Sekarang kita siap untuk teorema utama
bagian ini.

\medskip

\noindent}{\bf Teorema}
Misalkan $(X,\Sigma,\mu)$ dan $(Y,\Tau,\nu)$ adalah ruang ukur, dan misalkan
$\lambda$ adalah ukuran produk c.l.d.\ pada $X\times Y$. Misalkan $W$ suatu
ruang Banach dan $\phi:L^1(\mu)\times L^1(\nu)\to W$ suatu operator bilinear
terbatas. Maka terdapat tepat satu operator linear terbatas
$T:L^1(\lambda)\to W$ sedemikian sehingga $T(u\otimes v)=\phi(u,v)$ untuk setiap
$u\in L^1(\mu)$ dan $v\in L^1(\nu)$, serta $\|T\|=\|\phi\|$.

\proof{{\bf (a)} Inti argumen ini adalah fakta
berikut: jika $E_0,\ldots,E_n$ adalah himpunan terukur berukuran berhingga dalam
$X$, $F_0,\ldots,F_n$ adalah himpunan terukur berukuran berhingga dalam $Y$,
$a_0,\ldots,a_n\in\Bbb R$, dan
$\sum_{i=0}^na_i\chi(E_i\times F_i)=0\,\,\lambda$-hampir di mana-mana, maka
$\sum_{i=0}^na_i\phi(\chi E_i^{\ssbullet},\chi F_i^{\ssbullet})=0$ dalam $W$.
\Prf\ Kita
dapat menemukan keluarga saling lepas $\langle G_j\rangle_{j\le m}$ yang terdiri atas himpunan
terukur berukuran berhingga dalam $X$ sedemikian sehingga setiap $E_i$ dapat dinyatakan sebagai
gabungan suatu subkeluarga dari $G_j$; sehingga $\chi E_i$ dapat
dinyatakan dalam bentuk $\sum_{j=0}^mb_{ij}\chi G_j$ (lihat 122Ca).
Dengan cara yang sama, kita
dapat menemukan keluarga saling lepas $\langle H_k\rangle_{k\le l}$ yang terdiri atas himpunan
terukur berukuran berhingga dalam $Y$ sedemikian sehingga setiap $\chi F_i$ dapat dinyatakan
sebagai $\sum_{k=0}^lc_{ik}\chi H_k$. Sekarang

\Centerline{$\sum_{j=0}^m\sum_{k=0}^l
\bigl(\sum_{i=0}^na_ib_{ij}c_{ik}\bigr)
\chi(G_j\times H_k)
=\sum_{i=0}^na_i\chi(E_i\times F_i)=0\,\,\lambda\text{-a.e.}$}

\noindent Karena $G_j\times H_k$ saling lepas, dan
$\lambda(G_j\times H_k)=\mu G_j\cdot\nu H_k$ untuk setiap $j$, $k$,
diperoleh bahwa untuk
setiap $j\le m$, $k\le l$, berlaku salah satu dari $\mu G_j=0$, $\nu H_k=0$, atau
$\sum_{i=0}^na_ib_{ij}c_{ik}=0$. Dalam ketiga kasus ini,
$\sum_{i=0}^na_ib_{ij}c_{ik}\phi(\chi G_j^{\ssbullet},\chi
H_k^{\ssbullet})=0$ dalam $W$. Tetapi ini berarti bahwa

\Centerline{$0=\sum_{j=0}^m\sum_{k=0}^l
\bigl(\sum_{i=0}^na_ib_{ij}c_{ik}\bigr)
\phi(\chi G_j^{\ssbullet},\chi H_k^{\ssbullet})
=\sum_{i=0}^na_i\phi(\chi E_i^{\ssbullet},\chi F_i^{\ssbullet}),$}

\noindent sebagaimana dinyatakan.\ \Qed

\medskip

{\bf (b)} Dengan demikian, jika $E_0,\ldots,E_n$, $E'_0,\ldots,E'_m$ adalah
himpunan terukur berukuran berhingga dalam $X$, $F_0,\ldots,F_n$,
$F'_0,\ldots,F'_m$ adalah himpunan terukur berukuran berhingga dalam $Y$,
$a_0,\ldots,a_n,a'_0,\ldots,a'_m\in\Bbb R$, dan
$\sum_{i=0}^na_i\chi(E_i\times F_i)=\sum_{i=0}^ma'_i\chi(E'_i\times
F'_i)\,\,\lambda$-hampir di mana-mana, maka

\Centerline{$\sum_{i=0}^na_i\phi(\chi
E_i^{\ssbullet},\chi F_i^{\ssbullet})=\sum_{i=0}^ma'_i\phi(\chi
{E'_i}^{\ssbullet},\chi {F'_i}^{\ssbullet})$}

\noindent dalam $W$. Misalkan $M$ adalah subruang
linear dari $L^1(\lambda)$ yang direntang oleh

\Centerline{$\{\chi(E\times F)^{\ssbullet}:E\in\Sigma,\,\mu E<\infty,\,
  F\in\Tau,\,\nu F<\infty\};$}

\noindent maka terdapat tepat satu peta $T_0:M\to W$ sedemikian sehingga

\Centerline{$T_0(\sum_{i=0}^na_i\chi(E_i\times F_i)^{\ssbullet})
=\sum_{i=0}^na_i\phi(\chi E_i^{\ssbullet},\chi F_i^{\ssbullet})$}

\noindent kapan pun $E_0,\ldots,E_n$ adalah himpunan terukur berukuran
berhingga dalam $X$, $F_0,\ldots,F_n$ adalah himpunan terukur berukuran berhingga
dalam $Y$, dan $a_0,\ldots,a_n\in\Bbb R$. Tentu saja $T_0$ linear.

\medskip

{\bf (c)} Sebagian perhitungan yang sama menunjukkan bahwa
$\|T_0u\|\le\|\phi\|\|u\|_1$ untuk setiap $u\in M$. \Prf\ Jika $u\in M$,
maka, menurut argumen dalam (a), kita dapat menyatakan $u$ sebagai
$\sum_{j=0}^m\sum_{k=0}^la_{jk}\chi(G_j\times
H_k)^{\ssbullet}$, dengan $\langle G_j\rangle_{j\le m}$ dan $\langle
H_k\rangle_{k\le l}$ merupakan keluarga saling lepas dari himpunan berukuran berhingga.
Sekarang

$$\eqalign{\|T_0u\|
&=\|\sum_{j=0}^m\sum_{k=0}^la_{jk}\phi(\chi G_j^{\ssbullet},\chi
  H_k^{\ssbullet})\|
\le\sum_{j=0}^m\sum_{k=0}^l|a_{jk}|\|\phi(\chi G_j^{\ssbullet},\chi
  H_k^{\ssbullet})\|\cr
&\le\sum_{j=0}^m\sum_{k=0}^l|a_{jk}|\|\phi\|\|\chi
  G_j^{\ssbullet}\|_1\|\chi H_k^{\ssbullet}\|_1
=\|\phi\|\sum_{j=0}^m\sum_{k=0}^l|a_{jk}|\mu G_j\cdot\nu H_k\cr
&=\|\phi\|\sum_{j=0}^m\sum_{k=0}^l|a_{jk}|\lambda(G_j\times H_k)
=\|\phi\|\|u\|_1,\cr}$$

\noindent sebagaimana dinyatakan.\ \Qed

\medskip

{\bf (d)} Poin berikutnya ialah mengamati bahwa $M$ padat dalam
$L^1(\lambda)$ untuk $\|\,\|_1$. \Prf\ Dengan mengulangi gagasan di atas sekali
lagi, kita mengamati bahwa jika $E_0,\ldots,E_n$ adalah himpunan berukuran berhingga dalam
$X$ dan $F_0,\ldots,F_n$ adalah himpunan berukuran berhingga dalam $Y$, maka
$\chi(\bigcup_{i\le n}E_i\times F_i)^{\ssbullet}\in M$; ini
karena, dengan
menyatakan setiap $E_i$ sebagai gabungan $G_j$, dengan $G_j$ saling lepas,
kita mempunyai

\Centerline{$\bigcup_{i\le n}E_i\times F_i=\bigcup_{j\le m}G_j\times
F'_j,$}

\noindent dengan $F'_j=\bigcup\{F_i:G_j\subseteq E_i\}$ untuk setiap $j$;
sekarang $\langle G_j\times F'_j\rangle_{j\le m}$ saling lepas, sehingga

\Centerline{$\chi(\bigcup_{j\le m}G_j\times F'_j)^{\ssbullet}
=\sum_{j=0}^m\chi(G_j\times F'_j)^{\ssbullet}\in M.$}

\noindent Jadi 251Ie memberi tahu kita bahwa kapan pun $\lambda H<\infty$ dan
$\epsilon>0$, terdapat $G$ sedemikian sehingga $\lambda(H\symmdiff G)\le\epsilon$
dan $\chi G^{\ssbullet}\in M$; sekarang

\Centerline{$\|\chi H^{\ssbullet}-\chi
G^{\ssbullet}\|_1=\lambda(G\symmdiff H)\le\epsilon$,}

\noindent sehingga $\chi H^{\ssbullet}$
dapat dihampiri secara sebarang dekat oleh anggota $M$, dan termasuk dalam
penutup $\overline{M}$ dari $M$ dalam $L^1(\lambda)$. Karena $M$ adalah
subruang linear dari $L^1(\lambda)$, demikian pula $\overline{M}$ (2A4Cb);
oleh karena itu $\overline{M}$ memuat kelas ekuivalensi dari semua
fungsi sederhana-$\lambda$; tetapi semua ini padat dalam $L^1(\lambda)$
(242Mb), sehingga $\overline{M}=L^1(\lambda)$, sebagaimana dinyatakan.\ \Qed

\medskip

{\bf (e)} Karena $W$ adalah ruang Banach, diperoleh suatu
operator linear terbatas $T:L^1(\lambda)\to W$ yang memperluas $T_0$, dengan
$\|T\|=\|T_0\|\le\|\phi\|$ (2A4I). Sekarang $T(u\otimes v)=\phi(u,v)$ untuk
setiap $u\in L^1(\mu)$, $v\in L^1(\nu)$. \Prf\ Jika $u=\chi E^{\ssbullet}$ dan
$v=\chi F^{\ssbullet}$, dengan $E$, $F$ himpunan terukur berukuran
berhingga, maka

\Centerline{$T(u\otimes v)=T(\chi(E\times F)^{\ssbullet})
=T_0(\chi(E\times F)^{\ssbullet})
=\phi(\chi E^{\ssbullet},\chi F^{\ssbullet})=\phi(u,v).$}

\noindent Karena $\phi$ dan $\otimes$ bilinear dan $T$ linear,

\Centerline{$T(f^{\ssbullet}\otimes g^{\ssbullet})
=\phi(f^{\ssbullet},g^{\ssbullet})$}

\noindent kapan pun $f$ dan $g$ merupakan fungsi sederhana. Sekarang, kapan pun
$u\in L^1(\mu)$, $v\in L^1(\nu)$, dan $\epsilon>0$, terdapat fungsi sederhana
$f$, $g$ sedemikian sehingga $\|u-f^{\ssbullet}\|_1\le\epsilon$,
$\|v-g^{\ssbullet}\|_1\le\epsilon$ (242Mb lagi); sehingga

$$\eqalign{\|\phi(u,v)-\phi(f^{\ssbullet},g^{\ssbullet})\|
&\le\|\phi(u-f^{\ssbullet},v-g^{\ssbullet})\|
+\|\phi(u,g^{\ssbullet}-v)\|
+\|\phi(f^{\ssbullet}-u,v)\|\cr
&\le\|\phi\|(\epsilon^2+\epsilon\|u\|_1+\epsilon\|v\|_1).}$$

\noindent Demikian pula,

\Centerline{$\|u\otimes v-f^{\ssbullet}\otimes g^{\ssbullet}\|_1
\le\epsilon(\epsilon+\|u\|_1+\|v\|_1),$}

\noindent sehingga

\Centerline{$\|T(u\otimes v)-T(f^{\ssbullet}\otimes g^{\ssbullet})\|
\le\epsilon\|T\|(\epsilon+\|u\|_1+\|v\|_1);$}

\noindent karena $T(f^{\ssbullet}\otimes g^{\ssbullet})
=\phi(f^{\ssbullet},g^{\ssbullet})$,

\Centerline{$\|T(u\otimes v)-\phi(u,v)\|
\le\epsilon(\|T\|+\|\phi\|)(\epsilon+\|u\|_1
+\|v\|_1).$}

\noindent Karena $\epsilon$ sembarang, $T(u\otimes v)=\phi(u,v)$, sebagaimana
diperlukan.\ \Qed

\medskip

{\bf (f)} Argumen dalam (e) memastikan bahwa $\|T\|\le\|\phi\|$.
Karena $\|u\otimes v\|_1\le\|u\|_1\|v\|_1$ untuk setiap $u\in L^1(\mu)$ dan
$v\in L^1(\nu)$, $\|\phi(u,v)\|\le\|T\|\|u\|_1\|v\|_1$ untuk setiap $u$, $v$,
dan $\|\phi\|\le\|T\|$; jadi $\|T\|=\|\phi\|$.

\medskip

{\bf (g)} Jadi $T$ mempunyai sifat-sifat yang diperlukan. Untuk melihat bahwa operator ini
tunggal, kita hanya perlu mengamati bahwa setiap operator linear terbatas
$S:L^1(\lambda)\to W$ sedemikian sehingga $S(u\otimes v)=\phi(u,v)$ untuk setiap $u\in
L^1(\mu)$, $v\in L^1(\nu)$ harus sama dengan $T$ pada objek berbentuk
$\chi(E\times F)^{\ssbullet}$ dengan $E$ dan $F$ berukuran berhingga,
dan
karena itu pada setiap anggota $M$; karena $M$ padat dan $S$ maupun
$T$ kontinu, keduanya sama di seluruh $L^1(\lambda)$.
}


\leader{253G}{Struktur urutan pada \dvrocolon{$L^1$}}\cmmnt{ Dalam
253F saya memperlakukan ruang-ruang $L^1$ semata-mata sebagai
ruang linear bernorma. Namun, secara umum, struktur urutan suatu
ruang $L^1$ (lihat 242C) sama pentingnya dengan normanya. Pemetaan
$\otimes:L^1(\mu)\times L^1(\nu)\to L^1(\lambda)$ menghormati struktur
urutan ketiga ruang tersebut dalam pengertian kuat berikut.

\medskip

\noindent}{\bf Proposisi} Misalkan $(X,\Sigma,\mu)$ dan $(Y,\Tau,\nu)$
ruang-ruang ukur, dan $\lambda$ ukuran produk c.l.d.\ pada $X\times
Y$. Maka

(a) $u\otimes v\ge 0$ dalam
$L^1(\lambda)$ bilamana $u\ge 0$ dalam $L^1(\mu)$ dan $v\ge 0$ dalam
$L^1(\nu)$.

(b) Kerucut positif $\{w:w\ge 0\}$ dari
$L^1(\lambda)$ tepat merupakan selubung
konveks tertutup $C$ dari $\{u\otimes v:u\ge 0,\,v\ge 0\}$ dalam $L^1(\lambda)$.

*(c) Misalkan $W$ sebarang kisi Banach, dan $T:L^1(\lambda)\to W$
suatu operator linear terbatas. Maka
pernyataan-pernyataan berikut berekuivalen:

\quad(i) $Tw\ge 0$ dalam $W$ bilamana $w\ge 0$ dalam
$L^1(\lambda)$;

\quad(ii) $T(u\otimes v)\ge 0$ dalam $W$ bilamana $u\ge 0$ dalam
$L^1(\mu)$ dan $v\ge 0$ dalam $L^1(\nu)$.

\proof{{\bf (a)} Jika $u$, $v\ge 0$, maka keduanya dapat dinyatakan sebagai
$f^{\ssbullet}$, $g^{\ssbullet}$, dengan $f\in\eusm L^1(\mu)$,
$g\in\eusm L^1(\nu)$, $f\ge 0$ dan $g\ge 0$.
Karena $f\otimes g\ge 0$, maka $u\otimes v=(f\otimes g)^{\ssbullet}\ge 0$.

\medskip

{\bf (b)(i)} Tuliskan $L^1(\lambda)^+$ untuk $\{w:w\in L^1(\lambda),\,w\ge
0\}$. Maka $L^1(\lambda)^+$ merupakan himpunan konveks tertutup dalam $L^1(\lambda)$
(242De); menurut (a), himpunan itu memuat $u\otimes v$ bilamana $u\in L^1(\mu)^+$ dan
$v\in L^1(\nu)^+$, sehingga harus memuat $C$.

\medskip

\quad{\bf (ii)}($\alpha$)
Tentu saja $0=0\otimes 0\in C$. ($\beta$) Jika $u\in M$, sebagaimana didefinisikan dalam
pembuktian 253F, dan $u>0$, maka $u$ dapat dinyatakan sebagai $\sum_{j\le
m,k\le l}a_{jk}\chi(G_j\times H_k)^{\ssbullet}$, dengan $G_0,\ldots,G_m$
dan $H_0,\ldots,H_l$ merupakan barisan saling lepas dari himpunan-himpunan berukuran berhingga,
seperti dalam bagian (a) dari pembuktian 253F. Sekarang $a_{jk}$ hanya dapat negatif jika
$\chi(G_j\times H_k)^{\ssbullet}=0$, sehingga, dengan mengganti setiap $a_{jk}$ dengan
$\max(0,a_{jk})$ jika perlu, kita dapat mengandaikan bahwa $a_{jk}\ge 0$ untuk
semua $j$, $k$. Tidak semua $a_{jk}$ dapat bernilai nol, sehingga $a=\sum_{j\le
m,k\le l}a_{jk}>0$, dan

\Centerline{$u
=\sum_{j\le m,k\le l}
  \Bover{a_{jk}}a\cdot a\chi(G_j\times H_k)^{\ssbullet}
=\sum_{j\le m,k\le l}\Bover{a_{jk}}a
  \cdot(a\chi G_j^{\ssbullet})\otimes\chi H_k^{\ssbullet}
\in C$.}

\noindent ($\gamma$) Jika $w\in L^1(\lambda)^+$ dan $\epsilon>0$, nyatakan
$w$ sebagai $h^{\ssbullet}$, dengan $h\ge 0$ dalam $\eusm L^1(\lambda)$. Ada
suatu fungsi sederhana $h_1\ge 0$ sedemikian sehingga $h_1\leae h$ dan
$\int h\le\int h_1+\epsilon$. Nyatakan $h_1$ sebagai
$\sum_{i=0}^na_i\chi H_i$,
dengan $\lambda H_i<\infty$ dan $a_i\ge 0$ untuk setiap $i$, dan untuk setiap
$i\le n$ pilih himpunan-himpunan $G_{i0},\ldots,G_{im_i}\in\Sigma$,
$F_{i0},\ldots,F_{im_i}\in\Tau$, semuanya berukuran berhingga, sedemikian sehingga
$G_{i0},\ldots,G_{im_i}$ saling lepas dan
$\lambda(H_i\symmdiff\bigcup_{j\le m_i}G_{ij}\times
F_{ij})\le\epsilon/(n+1)(a_i+1)$, seperti dalam bagian (d) dari pembuktian 253F. Tetapkan

\Centerline{$w_0
=\sum_{i=0}^na_i\sum_{j=0}^{m_i}\chi(G_{ij}\times F_{ij})^{\ssbullet}$.}

\noindent Maka $w_0\in C$ karena $w_0\in M$ dan $w_0\ge 0$. Selain itu,

$$\eqalign{\|w-w_0\|_1
&\le\|w-h_1^{\ssbullet}\|_1+\|h_1^{\ssbullet}-w_0\|_1\cr
&\le\int(h-h_1)d\lambda
  +\sum_{i=0}^na_i\int|\chi H_i-\sum_{j=0}^{m_i}\chi(G_{ij}\times
  F_{ij})|d\lambda\cr
&\le\epsilon+\sum_{i=0}^na_i\lambda(H_i\symmdiff\bigcup_{j\le
m_i}G_{ij}\times F_{ij})
\le 2\epsilon.\cr}$$

\noindent Karena $\epsilon$ sebarang dan $C$ tertutup, $w\in C$. Karena
$w$ sebarang, $L^1(\lambda)^+\subseteq C$ dan $C=L^1(\lambda)^+$.

\medskip

{\bf (c)} Bagian (a) menyatakan bahwa (i)$\Rightarrow$(ii). Untuk implikasi
sebaliknya, kita memerlukan suatu bagian dari teori kisi Banach:
$W^+=\{w:w\in W,\,w\ge 0\}$ merupakan himpunan tertutup dalam $W$.
\Prf\ Jika $w$, $w'\in W$, maka

\Centerline{$w=(w-w')+w'\le|w-w'|+w'\le|w-w'|+|w'|$,}

\Centerline{$-w=(w'-w)-w'\le|w-w'|-w'\le|w-w'|+|w'|$,}

\Centerline{$|w|\le|w-w'|+|w'|$,
\quad$|w|-|w'|\le|w-w'|$,}

\noindent karena $|w|=w\vee(-w)$ dan urutan pada $W$ bersifat
invarian terhadap translasi (241Ec). Demikian pula, $|w'|-|w|\le|w-w'|$ dan
$||w|-|w'||\le|w-w'|$, sehingga $\||w|-|w'|\|\le\|w-w'\|$, menurut definisi
kisi Banach (242G). Dengan menetapkan $\phi(w)=|w|-w$, kita melihat bahwa
$\|\phi(w)-\phi(w')\|\le 2\|w-w'\|$ untuk semua $w$, $w'\in W$, sehingga
$\phi$ kontinu.

Sekarang, karena urutan invarian terhadap perkalian dengan skalar
positif,

\Centerline{$w\ge 0\iff 2w\ge 0\iff w\ge-w\iff w=|w|\iff\phi(w)=0$,}

\noindent sehingga $W^+=\{w:\phi(w)=0\}$ tertutup.\ \Qed

Sekarang andaikan bahwa (ii) benar, dan tetapkan
$C_1=\{w:w\in L^1(\lambda),\,Tw\ge 0\}$. Maka $C_1$ memuat
$u\otimes v$ bilamana $u$, $v\ge 0$; tetapi himpunan itu juga
konveks, karena $T$ linear, dan tertutup, karena $T$
kontinu dan $C_1=T^{-1}[W^+]$. Menurut (b), $C_1$ memuat
$\{w:w\in L^1(\lambda),\,w\ge 0\}$, sebagaimana dituntut oleh (i).
}%end of proof of 253G

\leader{253H}{\dvrocolon{Ekspektasi} bersyarat}\cmmnt{ Gagasan-gagasan
dalam bagian ini dan bagian sebelumnya memberi kita beberapa
contoh ekspektasi bersyarat yang terpenting.

\medskip

\noindent}{\bf Teorema} Misalkan $(X,\Sigma,\mu)$ dan $(Y,\Tau,\nu)$
ruang-ruang probabilitas lengkap, dengan produk c.l.d.
$(X\times Y,\Lambda,\lambda)$. Tetapkan $\Lambda_1=\{E\times Y:E\in\Sigma\}$.
Maka $\Lambda_1$ merupakan subaljabar-$\sigma$ dari $\Lambda$. Diberikan suatu
fungsi yang dapat diintegralkan terhadap $\lambda$ dan bernilai real, yakni $f$, tetapkan

\Centerline{$g(x,y)=\int f(x,z)\nu(dz)$}

\noindent bilamana $x\in X$, $y\in Y$ dan integral itu terdefinisi dalam
$\Bbb R$. Maka $g$ merupakan ekspektasi bersyarat dari $f$ pada $\Lambda_1$.

\proof{ Kita tahu bahwa $\Lambda_1\subseteq\Lambda$, menurut 251Ia, dan
$\Lambda_1$ merupakan aljabar-$\sigma$ himpunan karena $\Sigma$ demikian.
Teorema Fubini (252B, 252C)
menyatakan bahwa $f_1(x)=\int f(x,z)\nu(dz)$ terdefinisi untuk hampir setiap
$x$, dan oleh karena itu $g=f_1\otimes\chi Y$ terdefinisi hampir di mana-mana
dalam $X\times Y$. $f_1$ terukur secara virtual terhadap $\mu$; karena $\mu$
lengkap, $f_1$ terukur terhadap $\Sigma$, sehingga $g$ terukur terhadap $\Lambda_1$
(karena $\{(x,y):g(x,y)\le\alpha\}=\{x:f_1(x)\le\alpha\}\times Y$ untuk
setiap $\alpha\in\Bbb R$). Akhirnya, jika $W\in\Lambda_1$, maka
$W=E\times Y$ untuk suatu $E\in\Sigma$, sehingga

$$\eqalignno{\int_Wg\,d\lambda
&=\int (f_1\otimes\chi Y)\times(\chi E\otimes\chi Y)d\lambda
=\int f_1\times\chi E\,d\mu\int\chi Y\,d\nu\cr
\noalign{\noindent (menurut 253D)}
&=\iint\chi E(x)f(x,y)\nu(dy)\mu(dx)
=\int f\times\chi(E\times Y)d\lambda\cr
\noalign{\noindent (menurut teorema Fubini)}
&=\int_Wf\,d\lambda.\cr}$$

\noindent Jadi $g$ merupakan ekspektasi bersyarat dari $f$.
}%end of proof of 253H

\leader{253I}{}\cmmnt{ Inilah saat yang tepat untuk mengemukakan suatu
hasil berguna tentang produk ukuran integral tak tentu.

\medskip

\noindent}{\bf Proposisi} Misalkan $(X,\Sigma,\mu)$ dan $(Y,\Tau,\nu)$ adalah
ruang-ruang ukur, serta $f\in\eusm L^0(\mu)$, $g\in\eusm L^0(\nu)$
fungsi-fungsi nonnegatif. Misalkan $\mu'$, $\nuprime$ merupakan ukuran
integral tak tentu yang bersesuaian\cmmnt{ (lihat \S234)}. Misalkan $\lambda$
merupakan produk c.l.d.\ dari $\mu$ dan $\nu$, serta $\lambda'$ merupakan ukuran
integral tak tentu yang didefinisikan dari $\lambda$ dan
$f\otimes g\in\eusm L^0(\lambda)$\cmmnt{ (253Cb)}. Maka $\lambda'$ merupakan
produk c.l.d.\ dari $\mu'$ dan $\nuprime$.

\proof{ Tuliskan $\theta$ untuk produk c.l.d.\ dari $\mu'$ dan $\nuprime$.

\medskip

{\bf (a)} Jika kita mengganti $\mu$ dengan pelengkapannya, kita tidak mengubah $\mu'$
(234Ke); pada saat yang sama, kita tidak mengubah $\lambda$, menurut 251T.
Hal yang sama berlaku untuk $\nu$. Jadi cukup membuktikan hasil tersebut dengan
asumsi bahwa $\mu$ dan $\nu$ lengkap; dalam hal ini $f$ dan
$g$ terukur serta mempunyai domain-domain yang terukur.

Tetapkan $F=\{x:x\in\dom f,\,f(x)>0\}$ dan $G=\{y:y\in\dom g,\,g(y)>0\}$, sehingga
$F\times G=\{w:w\in\dom(f\otimes g),\,(f\otimes g)(w)>0\}$. Maka
$F$ koterabaikan terhadap $\mu'$ dan $G$ koterabaikan terhadap $\nuprime$, sehingga
$F\times G$ koterabaikan terhadap $\theta$ maupun
$\lambda'$. Karena $\theta$ dan $\lambda'$ keduanya
lengkap (251Ic, 234I), cukup ditunjukkan bahwa ukuran-ukuran subruang
$\theta_{F\times G}$, $\lambda'_{F\times G}$ pada $F\times G$ adalah
sama. Namun, perhatikan bahwa $\theta_{F\times G}$ dapat diidentifikasi dengan
produk dari $\mu_F'$ dan $\nu_G'$, dengan $\mu_F'$ dan $\nu_G'$ merupakan
ukuran-ukuran subruang masing-masing pada $F$, $G$ (251Q(ii-$\alpha$)). Pada
saat yang sama, $\mu_F'$ merupakan ukuran integral tak tentu yang didefinisikan dari
ukuran subruang $\mu_F$ pada $F$ dan fungsi $f\restr F$,
$\nu_G'$ merupakan ukuran integral tak tentu yang didefinisikan dari ukuran subruang
$\nu_G$ pada $G$ dan $g\restr G$, serta $\lambda'_{F\times G}$
didefinisikan dari ukuran subruang $\lambda_{F\times G}$ dan
$(f\restr F)\otimes(g\restr G)$. Akhirnya, sekali lagi menurut 251Q,
$\lambda_{F\times G}$ merupakan produk dari $\mu_F$ dan $\nu_G$.

Arti semua ini ialah bahwa cukup menangani kasus
di mana $F=X$ dan $G=Y$, yaitu, $f$ dan $g$ terdefinisi di setiap titik dan
positif ketat; inilah yang saya andaikan mulai sekarang.

\medskip

{\bf (b)} Dalam hal ini $\dom\mu'=\Sigma$ dan $\dom\nuprime=\Tau$ (234La).
Demikian pula, $\dom\lambda'=\Lambda$ tidak lain adalah domain dari $\lambda$. Tetapkan

\Centerline{$F_n=\{x:x\in X,\,2^{-n}\le f(x)\le 2^n\}$,
\quad$G_n=\{y:y\in Y,\,2^{-n}\le g(y)\le 2^n\}$}

\noindent untuk $n\in\Bbb N$.

\medskip

{\bf (c)} Tetapkan

\Centerline{$\Cal A=\{W:W\in\dom\theta\cap\dom\lambda'$,
$\theta(W)=\lambda'(W)\}$.}

\noindent Jika $\mu'E$ dan $\nuprime H$ terdefinisi dan berhingga, maka
$f\times\chi E$ dan $g\times\chi H$ dapat diintegralkan, sehingga

$$\eqalign{\lambda'(E\times H)
&=\int(f\otimes g)\times\chi(E\times H)d\lambda
=\int(f\times\chi E)\otimes(g\times\chi H)d\lambda\cr
&=\int f\times\chi E\,d\mu\cdot\int g\times\chi H\,d\nu
=\theta(E\times H)\cr}$$

\noindent menurut 253D dan 251Ia, yaitu, $E\times H\in\Cal A$. Jika sekarang kita
meninjau
$\Cal A_{EH}=\{W:W\subseteq X\times Y$, $W\cap(E\times H)\in\Cal A\}$,
maka kita melihat bahwa

\inset{$\Cal A_{EH}$ memuat $E'\times H'$ untuk setiap $E'\in\Sigma$,
$H'\in\Tau$,}

\inset{jika $\sequencen{W_n}$ suatu barisan tak menurun dalam
$\Cal A_{EH}$, maka $\bigcup_{n\in\Bbb N}W_n\in\Cal A_{EH}$,}

\inset{jika $W$, $W'\in\Cal A_{EH}$ dan $W\subseteq W'$, maka
$W'\setminus W\in\Cal A_{EH}$.}

\noindent Dengan demikian, $\Cal A_{EH}$ merupakan kelas Dynkin dari subhimpunan-subhimpunan
$X\times Y$, dan menurut Teorema Kelas Monoton (136B), memuat
aljabar-$\sigma$ yang dibangkitkan oleh
$\{E'\times H':E'\in\Sigma,\,H'\in\Tau\}$, yaitu
$\Sigma\tensorhat\Tau$.

\medskip

{\bf (d)} Sekarang andaikan bahwa $W\in\Lambda$. Dalam hal ini
$W\in\dom\theta$ dan $\theta W\le\lambda'W$. \Prf\ Ambil $n\in\Bbb N$,
serta $E\in\Sigma$, $H\in\Tau$ sedemikian sehingga $\mu'E$ dan $\nuprime H$
keduanya berhingga. Tetapkan $E'=E\cap F_n$, $H'=H\cap G_n$ dan $W'=W\cap(E'\times
H')$. Maka $W'\in\Lambda$, sedangkan $\mu E'\le 2^n\mu'E$ dan $\nu H'\le
2^n\nuprime H$ berhingga. Menurut 251Ib, terdapat $V\in\Sigma\tensorhat\Tau$
sedemikian sehingga $V\subseteq W'$ dan $\lambda V=\lambda W'$. Demikian pula, terdapat
$V'\in\Sigma\tensorhat\Tau$ sedemikian sehingga $V'\subseteq(E'\times
H')\setminus W'$ dan $\lambda V'=\lambda((E'\times H')\setminus W')$.
Ini berarti bahwa $\lambda((E'\times H')\setminus(V\cup V'))=0$, sehingga
$\lambda'((E'\times H')\setminus(V\cup V'))=0$. Namun,
$(E'\times H')\setminus(V\cup V')\in\Cal A$, menurut (c), sehingga $\theta((E'\times H')\setminus(V\cup V'))=0$ dan $W'\in\dom\theta$, sedangkan

\Centerline{$\theta W'=\theta V=\lambda'V
\le\lambda'W$.}

Karena $E$ dan $H$ sebarang, $W\cap(F_n\times G_n)\in\dom\theta$
(251H) dan $\theta(W\cap(F_n\times G_n))\le\lambda'W$. Karena
$\sequencen{F_n}$, $\sequencen{G_n}$ merupakan barisan-barisan tak menurun dengan
gabungan masing-masing $X$, $Y$,

\Centerline{$\theta W=\sup_{n\in\Bbb N}\theta(W\cap(F_n\times G_n))
\le\lambda'W$. \Qed}

\medskip

{\bf (e)} Dengan cara yang sama, $\lambda'W$ terdefinisi dan tidak lebih besar
daripada $\theta W$ untuk setiap $W\in\dom\theta$. \Prf\ Argumen-argumennya sangat
mirip, tetapi tampaknya diperlukan suatu penyempurnaan pada tahap terakhir.
Ambil $n\in\Bbb N$, serta $E\in\Sigma$, $H\in\Tau$ sedemikian sehingga $\mu E$
dan $\nu H$ keduanya berhingga. Tetapkan $E'=E\cap F_n$, $H'=H\cap G_n$ dan
$W'=W\cap(E'\times H')$. Maka $W'\in\dom\theta$, sedangkan $\mu'E'\le
2^n\mu E$ dan $\nuprime H'\le 2^n\nu H$ berhingga. Kali ini, terdapat
$V$, $V'\in\Sigma\tensorhat\Tau$ sedemikian sehingga $V\subseteq W'$,
$V'\subseteq(E'\times H')\setminus W'$, $\theta V=\theta W'$ dan
$\theta V'=\theta((E'\times H')\setminus W')$. Dengan demikian,

\Centerline{$\lambda'V+\lambda'V'=\theta V+\theta V'
=\theta(E'\times H')=\lambda'(E'\times H')$,}

\noindent sehingga $\lambda'W'$ terdefinisi dan sama dengan $\theta W'$.

Ini berarti bahwa $W\cap(F_n\times G_n)\cap(E\times H)\in\Cal A$
bilamana $\mu E$ dan $\nu H$ berhingga. Jadi
$W\cap(F_n\times G_n)\in\Lambda$, menurut 251H;
karena $n$ sebarang, $W\in\Lambda$ dan $\lambda'W$ terdefinisi.

\Quer\ Andaikan, jika mungkin, bahwa $\lambda'W>\theta W$. Maka terdapat
suatu $n\in\Bbb N$ sedemikian sehingga $\lambda'(W\cap(F_n\times G_n))>\theta W$.
Karena $\lambda$ semihingga, 213B menyatakan bahwa terdapat suatu
fungsi sederhana-$\lambda$ $h$ sedemikian sehingga
$h\le(f\otimes g)\times\chi(W\cap(F_n\times G_n))$ dan
$\int h\,d\lambda>\theta W$; dengan menetapkan $V=\{(x,y):h(x,y)>0\}$, kita melihat bahwa $V\subseteq W\cap(F_n\times G_n)$,
$\lambda V$ terdefinisi dan berhingga, serta $\lambda'V>\theta W$. Sekarang pasti
terdapat himpunan-himpunan $E\in\Sigma$, $H\in\Tau$ sedemikian sehingga $\mu E$ dan $\nu H$
keduanya berhingga dan
$\lambda(V\setminus(E\times H))<4^{-n}(\lambda'V-\theta W)$. Namun, dalam hal ini $V\in\Lambda\subseteq\dom\theta$ (menurut (d)), sehingga kita dapat menerapkan argumen tepat di atas pada $V$ dan menyimpulkan bahwa
$V\cap(E\times H)=V\cap(F_n\times G_n)\cap(E\times H)$ termasuk dalam
$\Cal A$. Sekarang,

$$\eqalign{\lambda'V
&=\lambda'(V\cap(E\times H))+\lambda'(V\setminus(E\times H))\cr
&\le\theta(V\cap(E\times H))+4^n\lambda(V\setminus(E\times H))
<\theta V+\lambda'V-\theta W
\le\lambda'V,\cr}$$

\noindent suatu kemustahilan.\ \Bang

Jadi $\lambda'W$ terdefinisi dan tidak lebih besar daripada $\theta W$.\
\Qed

\medskip

{\bf (f)} Dengan menggabungkan ini dengan (d), kita melihat bahwa $\lambda'=\theta$,
sebagaimana diklaim.
}%end of proof of 253I

\cmmnt{\medskip

\noindent{\bf Catatan} Jika $\mu'$ dan $\nuprime$ berhingga total, sehingga
keduanya ‘benar-benar kontinu’ terhadap $\mu$ dan $\nu$ dalam pengertian
232Ab, maka $f$ dan $g$ dapat diintegralkan, sehingga $f\otimes g$ dapat
diintegralkan terhadap $\lambda$, dan $\theta=\lambda'$ benar-benar kontinu terhadap
$\lambda$.

Pembuktian di atas dapat disederhanakan dengan menggunakan bagian-bagian dari teori umum
ruang lengkap yang ditentukan secara lokal, yang akan diberikan dalam \S412 %412J
pada Jilid 4.}%end of comment


\leader{*253J}{Integral \dvrocolon{atas}}\cmmnt{ Gagasan 253D
dapat diulangi dalam istilah integral atas, sebagai berikut.

\medskip

\noindent}{\bf Proposisi} Misalkan $(X,\Sigma,\mu)$ dan $(Y,\Tau,\nu)$
ruang-ruang ukur $\sigma$-hingga, dengan ukuran produk c.l.d.\ $\lambda$.
Maka untuk sebarang fungsi $f$ dan $g$, yang masing-masing didefinisikan
pada subhimpunan koterabaikan dari $X$ dan $Y$, dan bernilai dalam
$[0,\infty]$,

\Centerline{$\overlineint f\otimes g\,d\lambda
=\overlineint fd\mu\cdot\overlineint g\,d\nu$.}

\cmmnt{\medskip

\noindent{\bf Catatan} Di sini $(f\otimes g)(x,y)=f(x)g(y)$ untuk semua
$x\in\dom f$ dan $y\in\dom g$, dengan menetapkan $0\cdot\infty=0$,
seperti dalam \S135.
}%end of comment

\proof{{\bf (a)} Mula-mula saya tunjukkan bahwa
$\overline{\int}f\otimes g\le\overline{\int}f\overline{\int}g$.   \Prf\
Jika $\overline{\int}f=0$, maka $f=0$ hampir di mana-mana, sehingga
$f\otimes g=0$ hampir di mana-mana.\ Dengan demikian, hasilnya langsung diperoleh.
Argumen yang sama berlaku jika $\overline{\int}g=0$. Jika
$\overline{\int}f$ dan $\overline{\int}g$ keduanya tidak nol, dan salah
satunya tak hingga, hasilnya trivial. Jadi, andaikan keduanya
berhingga. Dalam hal ini terdapat $f_0$, $g_0$ terintegralkan sedemikian
sehingga $f\leae f_0$, $g\leae g_0$,
$\overline{\int}f=\int f_0$ dan $\overline{\int}g=\int g_0$ (133Ja/135Ha).
Jadi $f\otimes g\leae f_0\otimes g_0$, dan

\Centerline{$
\overlineint f\otimes g\le\int f_0\otimes g_0
=\int f_0\int g_0=\overlineint f\overlineint g$,}

\noindent menurut 253D.\ \Qed

\medskip

{\bf (b)} Untuk ketaksamaan sebaliknya, kita hanya perlu meninjau kasus
ketika $\overline{\int}f\otimes g$ berhingga, sehingga terdapat suatu
fungsi yang terintegralkan terhadap $\lambda$, yaitu $h$, sedemikian sehingga
$f\otimes g\leae h$ dan
$\overline{\int}f\otimes g=\int h$. Tetapkan

\Centerline{$f_0(x)=\int h(x,y)\nu(dy)$}

\noindent kapan pun nilai ini terdefinisi dalam $\Bbb R$, yaitu hampir
di mana-mana, menurut teorema Fubini (252B-252C). Maka
$f_0(x)\ge f(x)\overline{\int}g\,d\nu$ untuk setiap
$x\in\dom f_0\cap\dom f$, yang merupakan himpunan koterabaikan dalam $X$; jadi

\Centerline{$
\overlineint f\otimes g=\int h\,d\lambda
=\int f_0d\mu\ge\overlineint f\overlineint g$,}

\noindent sebagaimana diperlukan.
}%end of proof of 253J

\leader{*253K}{}\cmmnt{ Argumen serupa berlaku untuk integral atas
dari jumlah, sebagai berikut.

\medskip

\noindent}{\bf Proposisi} Misalkan $(X,\Sigma,\mu)$ dan $(Y,\Tau,\nu)$
ruang-ruang probabilitas, dengan ukuran produk c.l.d.\ $\lambda$. Maka
untuk sebarang fungsi bernilai real $f$, $g$ yang masing-masing didefinisikan
pada subhimpunan koterabaikan dari $X$, $Y$,

\Centerline{$
\overlineint f(x)+g(y)\,\lambda(d(x,y))
=\overlineint f(x)\mu(dx)+\overlineint g(y)\nu(dy)$,}

\noindent setidaknya ketika ruas kanan terdefinisi dalam
$[-\infty,\infty]$.

\proof{Tetapkan $h(x,y)=f(x)+g(y)$ untuk $x\in\dom f$ dan
$y\in\dom g$, sehingga $\dom h$ koterabaikan terhadap $\lambda$.

\medskip

{\bf (a)} Seperti dalam 253J, saya mulai dengan menunjukkan bahwa
$\overline{\int}h\le\overline{\int}f+\overline{\int}g$.   \Prf\ Jika
salah satu dari $\overline{\int}f$ atau $\overline{\int}g$ adalah
$\infty$, pernyataan ini trivial. Jika tidak, ambillah fungsi-fungsi
terintegralkan $f_0$, $g_0$ sedemikian sehingga
$f\leae f_0$ dan $g\leae g_0$. Tetapkan
$h_0=(f_0\otimes\chi Y)+(\chi X\otimes g_0)$; maka
$h\le h_0\,\,\lambda$-hampir di mana-mana, sehingga

\Centerline{$
\overlineint h\,d\lambda\le\int h_0d\lambda
=\int f_0d\mu+\int g_0d\nu$.}

\noindent Karena $f_0$, $g_0$ sebarang,
$\overline{\int}h\le\overline{\int}f+\overline{\int}g$.\ \Qed

\medskip

{\bf (b)} Untuk ketaksamaan sebaliknya, andaikan bahwa $h\le h_0$ untuk
$\lambda$-hampir setiap $(x,y)$, dengan $h_0$ terintegralkan terhadap
$\lambda$. Tetapkan $f_0(x)=\int h_0(x,y)\nu(dy)$ kapan pun nilai ini
terdefinisi dalam $\Bbb R$. Maka
$f_0(x)\ge f(x)+\overline{\int}g\,d\nu$ kapan pun $x\in\dom
f\cap\dom f_0$, sehingga

\Centerline{$\int h_0\,d\lambda
=\int f_0d\mu\ge\overlineint fd\mu+\overlineint g\,d\nu$.}

\noindent Karena $h_0$ sebarang,
$\overline{\int}h\ge\overline{\int}f+\overline{\int}g$, sebagaimana diperlukan.
}%end of proof of 253K

\leader{253L}{Ruang kompleks} Seperti biasa, gagasan-gagasan\cmmnt{ dari 253F
dan 253H} berlaku pada dasarnya tanpa perubahan untuk ruang $L^1$
kompleks. Dengan menuliskan $L^1_{\Bbb C}(\mu)$, dan seterusnya, untuk
ruang-ruang $L^1$ kompleks yang bersangkutan, kita memperoleh
hasil-hasil\cmmnt{ berikut}.
Di seluruh bagian ini, misalkan $(X,\Sigma,\mu)$ dan $(Y,\Tau,\nu)$
ruang-ruang ukur, dan $\lambda$ ukuran produk c.l.d.\ pada
$X\times Y$.

\spheader 253La Jika $f\in\eusm L^0_{\Bbb C}(\mu)$ dan
$g\in\eusm L^0_{\Bbb C}(\nu)$ maka $f\otimes g$, yang didefinisikan oleh rumus
$(f\otimes g)(x,y)=f(x)g(y)$ untuk $x\in\dom f$ dan $y\in\dom g$, termasuk
dalam $\eusm L^0_{\Bbb C}(\lambda)$.

\spheader 253Lb Jika $f\in\eusm L^1_{\Bbb C}(\mu)$ dan
$g\in\eusm L^1_{\Bbb C}(\nu)$ maka
$f\otimes g\in\eusm L^1_{\Bbb C}(\lambda)$ dan
$\int f\otimes g\,d\lambda=\int fd\mu\int g\,d\nu$.

\spheader 253Lc Kita mempunyai operator bilinear $(u,v)\mapsto u\otimes v:
L^1_{\Bbb C}(\mu)\times L^1_{\Bbb C}(\nu)\to L^1_{\Bbb C}(\lambda)$
yang didefinisikan dengan menuliskan $f^{\ssbullet}\otimes g^{\ssbullet}=(f\otimes
g)^{\ssbullet}$ untuk semua
$f\in\eusm L^1_{\Bbb C}(\mu)$, $g\in \eusm
L^1_{\Bbb C}(\nu)$.

\spheader 253Ld Jika $W$ sebarang ruang Banach kompleks dan
$\phi:L^1_{\Bbb C}(\mu)\times L^1_{\Bbb C}(\nu)\to W$ sebarang operator
bilinear terbatas, maka terdapat suatu operator linear terbatas yang unik
$T:L^1_{\Bbb C}(\lambda)\to W$ sedemikian sehingga
$T(u\otimes v)=\phi(u,v)$ untuk setiap $u\in L^1_{\Bbb C}(\mu)$ dan
$v\in L^1_{\Bbb C}(\nu)$, dan $\|T\|=\|\phi\|$.

\spheader 253Le Jika $\mu$ dan $\nu$ merupakan ukuran-ukuran probabilitas
lengkap, dan $\Lambda_1=\{E\times Y:E\in\Sigma\}$, maka untuk sebarang
$f\in\eusm
L^1_{\Bbb C}(\lambda)$ kita mempunyai ekspektasi bersyarat
$g$ dari $f$ pada $\Lambda_1$ yang diberikan dengan menetapkan
$g(x,y)=\int f(x,z)\nu(dz)$ kapan pun nilai ini terdefinisi.

\exercises{
\leader{253X}{Latihan dasar $\pmb{>}$(a)}
%\spheader 253Xa
Misalkan $U$, $V$ dan $W$ ruang-ruang linear. Tunjukkan bahwa himpunan
operator bilinear dari $U\times V$ ke $W$ mempunyai struktur linear alami
yang bersesuaian dengan struktur pada $\eurm L(U;\eurm L(V;W))$ dan
$\eurm L(V;\eurm
L(U;W))$, dengan menuliskan $\eurm L(U;W)$
untuk ruang linear semua operator linear dari $U$ ke $W$.

\sqheader 253Xb Misalkan $U$, $V$ dan $W$ ruang-ruang bernorma. (i)
Tunjukkan bahwa untuk suatu operator bilinear $\phi:U\times V\to W$
pernyataan-pernyataan berikut ekuivalen:
($\alpha$) $\phi$ terbatas dalam arti 253Ab; ($\beta$) $\phi$
kontinu; ($\gamma$) $\phi$ kontinu di suatu titik pada $U\times
V$. (ii) Tunjukkan bahwa ruang operator bilinear terbatas dari $U\times V$
ke $W$ merupakan subruang linear dari ruang semua operator bilinear dari
$U\times V$ ke $W$, dan bahwa fungsional $\|\,\|$ yang didefinisikan
dalam 253Ab merupakan suatu norma, yang bersesuaian dengan norma-norma
$\eurm B(U;\eurm B(V;W))$ dan $\eurm
B(V;\eurm B(U;W))$, dengan menuliskan $\eurm B(U;W)$ untuk ruang bernorma
semua operator linear terbatas dari $U$ ke $W$.

\spheader 253Xc Misalkan
$(X_1,\Sigma_1,\mu_1),\ldots,(X_n,\Sigma_n,\mu_n)$
ruang-ruang ukur, dan $\lambda$ ukuran produk c.l.d.\ pada
$X_1\times\ldots\times X_n$, sebagaimana dideskripsikan dalam 251W.
Misalkan $W$ ruang Banach, dan andaikan bahwa
$\phi:L^1(\mu_1)\times\ldots\times L^1(\mu_n)\to W$ bersifat
{\bf multilinear} (yaitu linear dalam setiap variabel
secara terpisah) dan {\bf terbatas} (yaitu,
$\|\phi\|=\sup\{\|\phi(u_1,\ldots,u_n)\|:
\|u_i\|_1\le 1\Forall i\le n\}<\infty$). Tunjukkan bahwa terdapat suatu
operator linear terbatas yang unik
$T:L^1(\lambda)\to W$ sedemikian sehingga $T\otimes=\phi$, dengan
$\otimes:L^1(\mu_1)\times\ldots\times L^1(\mu_n)\to L^1(\lambda)$
merupakan suatu operator multilinear kanonik (yang harus didefinisikan).

\spheader 253Xd Misalkan $(X,\Sigma,\mu)$ dan $(Y,\Tau,\nu)$ ruang-ruang
ukur, dan $\lambda$ ukuran produk c.l.d.\ pada $X\times Y$. Tunjukkan
bahwa jika $A\subseteq L^1(\mu)$ dan $B\subseteq L^1(\nu)$ keduanya
terintegralkan secara seragam, maka $\{u\otimes v:u\in A,\,v\in B\}$
terintegralkan secara seragam dalam $L^1(\lambda)$.

\sqheader 253Xe Misalkan $(X,\Sigma,\mu)$ dan $(Y,\Tau,\nu)$ ruang-ruang
ukur dan $\lambda$ ukuran produk c.l.d.\ pada $X\times Y$.
Tunjukkan bahwa

\quad (i) kita mempunyai operator bilinear
$(u,v)\mapsto u\otimes v:L^0(\mu)\times L^0(\nu)\to L^0(\lambda)$
yang diberikan dengan menetapkan
$f^{\ssbullet}\otimes g^{\ssbullet}=(f\otimes g)^{\ssbullet}$ untuk semua
$f\in\eusm L^0(\mu)$ dan $g\in \eusm L^0(\nu)$;

\quad (ii) jika $1\le p\le\infty$ maka $u\otimes v\in L^p(\lambda)$ dan
$\|u\otimes v\|_p=\|u\|_p\|v\|_p$ untuk semua $u\in L^p(\mu)$ dan
$v\in L^p(\nu)$;

\quad (iii) jika $u$, $u'\in L^2(\mu)$ dan $v$, $v'\in L^2(\nu)$ maka
hasil kali dalam $\innerprod{u\otimes v}{u'\otimes v'}$, yang diambil
dalam $L^2(\lambda)$, sama dengan
$\innerprod{u}{u'}\innerprod{v}{v'}$;

\quad (iv) pemetaan
$(u,v)\mapsto u\otimes v:L^0(\mu)\times L^0(\nu)\to L^0(\lambda)$
kontinu jika $L^0(\mu)$, $L^0(\nu)$ dan $L^0(\lambda)$
semuanya dilengkapi dengan topologi konvergensi dalam ukuran masing-masing.

\spheader 253Xf Dalam 253Xe, andaikan bahwa $\mu$ dan $\nu$ semihingga.
Tunjukkan bahwa jika $u_0,\ldots,u_n$ merupakan anggota-anggota yang bebas
linear dari $L^0(\mu)$ dan $v_0,\ldots,v_n\in L^0(\nu)$ tidak semuanya
$0$, maka $\sum_{i=0}^nu_i\otimes v_i\ne 0$ dalam $L^0(\lambda)$.
({\it Petunjuk\/}: mulailah dengan mencari himpunan-himpunan
$E\in\Sigma$, $F\in\Tau$ berukuran berhingga sedemikian sehingga
$u_0\times\chi E^{\ssbullet},\ldots,u_n\times\chi E^{\ssbullet}$
bebas linear dan
$v_0\times\chi F^{\ssbullet},\ldots,v_n\times\chi F^{\ssbullet}$
tidak semuanya $0$.)

\spheader 253Xg Dalam 253Xe, andaikan bahwa $\mu$ dan $\nu$ semihingga.
Jika $U$, $V$ masing-masing merupakan subruang linear dari $L^0(\mu)$
dan $L^0(\nu)$, tuliskan $U\otimes V$ untuk subruang linear dari
$L^0(\lambda)$ yang dibangkitkan oleh
$\{u\otimes v:u\in U,\,v\in V\}$. Tunjukkan bahwa jika $W$ sebarang ruang
linear dan $\phi:U\times V\to W$ suatu operator bilinear, terdapat suatu
operator linear unik $T:U\otimes V\to W$
sedemikian sehingga $T(u\otimes v)=\phi(u,v)$ untuk semua $u\in U$,
$v\in V$. ({\it Petunjuk\/}: mulailah dengan menunjukkan bahwa jika
$u_0,\ldots,u_n\in U$ dan $v_0,\ldots,v_n\in V$ sedemikian sehingga
$\sum_{i=0}^nu_i\otimes v_i=0$, maka
$\sum_{i=0}^n\phi(u_i,v_i)=0$ -- lakukan ini dengan menyatakan $u_i$
sebagai kombinasi linear dari suatu keluarga bebas linear, lalu terapkan
253Xf.)

\sqheader 253Xh Misalkan $(X,\Sigma,\mu)$ dan $(Y,\Tau,\nu)$ ruang-ruang
probabilitas lengkap, dengan ukuran produk c.l.d.\ $\lambda$. Andaikan
bahwa $p\in[1,\infty]$ dan bahwa $f\in\eusm L^p(\lambda)$. Tetapkan
$g(x)=\int f(x,y)\nu(dy)$ kapan pun nilai ini terdefinisi. Tunjukkan bahwa
$g\in\eusm L^p(\mu)$ dan bahwa $\|g\|_p\le\|f\|_p$.
({\it Petunjuk\/}: 253H, 244M.)

\spheader 253Xi Misalkan $(X,\Sigma,\mu)$ dan $(Y,\Tau,\nu)$ ruang-ruang
ukur, dengan ukuran produk c.l.d.\ $\lambda$, dan
$p\in\coint{1,\infty}$. Tunjukkan bahwa
$\{w:w\in L^p(\lambda),\,w\ge 0\}$ merupakan selubung konveks tertutup
dalam $L^p(\lambda)$ dari
$\{u\otimes v:u\in L^p(\mu),\,v\in L^p(\nu),\,u\ge 0,\,v\ge 0\}$
(lihat 253Xe(ii) di atas).

\leader{253Y}{Latihan lebih lanjut (a)}
%\spheader 253Ya
Misalkan $(X,\Sigma,\mu)$ dan $(Y,\Tau,\nu)$ ruang-ruang ukur, dan
$\lambda_0$ ukuran produk primitif pada $X\times Y$. Tunjukkan bahwa jika
$f\in\eusm L^0(\mu)$ dan $g\in\eusm L^0(\nu)$, maka
$f\otimes g\in\eusm L^0(\lambda_0)$.
%253C

\spheader 253Yb Misalkan $(X,\Sigma,\mu)$ dan
$(Y,\Tau,\nu)$ ruang-ruang ukur, dan $\lambda_0$ ukuran produk
primitif pada $X\times Y$. Tunjukkan bahwa jika $f\in\eusm L^1(\mu)$ dan
$g\in\eusm L^1(\nu)$, maka $f\otimes g\in\eusm L^1(\lambda_0)$ dan
$\int f\otimes g\,d\lambda_0=\int f\,d\mu\int g\,d\nu$.
%253Ya, 253C

\spheader 253Yc Misalkan $(X,\Sigma,\mu)$ dan $(Y,\Tau,\nu)$
ruang-ruang ukur, dan $\lambda_0$, $\lambda$ masing-masing ukuran produk
primitif dan c.l.d.\
pada $X\times Y$. Tunjukkan bahwa pembenaman
$\eusm L^1(\lambda_0)\embedsinto\eusm L^1(\lambda)$
menginduksi isomorfisme kisi Banach antara $L^1(\lambda_0)$ dan
$L^1(\lambda)$.
%253D

\spheader 253Yd Misalkan $(X,\Sigma,\mu)$, $(Y,\Tau,\nu)$ ruang-ruang
ukur yang dapat dilokalkan secara ketat, dengan ukuran produk c.l.d.\ $\lambda$.
Tunjukkan bahwa $L^{\infty}(\lambda)$ dapat diidentifikasi
dengan $L^1(\lambda)^*$. Tunjukkan bahwa di bawah identifikasi ini
$\{w:w\in L^{\infty}(\lambda),\,w\ge 0\}$ merupakan selubung konveks yang
tertutup dalam topologi lemah* dari
$\{u\otimes v:u\in L^{\infty}(\mu),\,v\in L^{\infty}(\nu),
\,u\ge 0,\,v\ge 0\}$.
%253G, 253Xe

\spheader 253Ye Carilah suatu versi 253J yang berlaku ketika salah satu
dari $\mu$, $\nu$ tidak $\sigma$-hingga.
%253J

\spheader 253Yf Misalkan $(X,\Sigma,\mu)$ sebarang ruang ukur dan $V$
sebarang ruang Banach. Tuliskan $\eusm L^1_V=\eusm L^1_V(\mu)$ untuk
himpunan fungsi-fungsi $f$
sedemikian sehingga ($\alpha$) $\dom f$ merupakan subhimpunan koterabaikan
dari $X$ ($\beta$) $f$ bernilai dalam $V$ ($\gamma$) terdapat suatu
himpunan koterabaikan
$D\subseteq\dom
f$ sedemikian sehingga $f[D]$ separabel dan $D\cap f^{-1}[G]\in\Sigma$
untuk setiap himpunan terbuka $G\subseteq V$ ($\delta$) integral
$\int\|f(x)\|\mu(dx)$ berhingga. (Ini adalah fungsi-fungsi {\bf yang
terintegralkan menurut Bochner} dari $X$ ke
$V$.) Untuk $f$, $g\in\eusm L^1_V$ tuliskan $f\sim g$ jika
$f=g\,\,\mu$-hampir di mana-mana; misalkan $L^1_V$ himpunan kelas-kelas
ekuivalensi dalam $\eusm L^1_V$ terhadap $\sim$. Tunjukkan bahwa

\quad (i) $f+g$, $cf\in\eusm L^1_V$ untuk semua $f$,
$g\in\eusm L^1_V$, $c\in\Bbb R$;

\quad (ii) $L^1_V$ mempunyai struktur ruang linear alami, yang didefinisikan
dengan menuliskan $f^{\ssbullet}+g^{\ssbullet}=(f+g)^{\ssbullet}$,
$cf^{\ssbullet}=(cf)^{\ssbullet}$ untuk $f$, $g\in\eusm L^1_V$ dan
$c\in\Bbb R$;

\quad (iii) $L^1_V$ mempunyai norma $\|\,\|$, yang didefinisikan dengan
menuliskan $\|f^{\ssbullet}\|=\int\|f(x)\|\mu(dx)$ untuk
$f\in\eusm L^1_V$;

\quad (iv) $L^1_V$ merupakan ruang Banach dengan norma ini;

\quad (v) terdapat suatu pemetaan alami
$\otimes:\eusm L^1\times V\to\eusm L^1_V$ yang didefinisikan dengan
menuliskan $(f\otimes v)(x)=f(x)v$ ketika
$f\in\eusm L^1=\eusm L^1_{\Bbb R}(\mu)$,
$v\in V$ dan $x\in\dom f$;

\quad (vi) terdapat suatu operator bilinear kanonik
$\otimes:L^1\times V\to L^1_V$ yang didefinisikan dengan menuliskan
$f^{\ssbullet}\otimes v=(f\otimes v)^{\ssbullet}$ untuk
$f\in\eusm L^1$ dan
$v\in V$;

\quad (vii) kapan pun $W$ suatu ruang Banach dan
$\phi:L^1\times V\to W$ suatu operator bilinear terbatas,
terdapat suatu operator linear terbatas yang unik
$T:L^1_V\to W$ sedemikian sehingga $T(u\otimes v)=\phi(u,v)$
untuk semua $u\in L^1$ dan $v\in V$, dan $\|T\|=\|\phi\|$.
%new 2008
(Ketika $W=V$ dan
$\phi(u,v)=(\int u)v$ untuk $u\in L^1$ dan $v\in V$, $Tf^{\ssbullet}$
disebut {\bf integral Bochner} dari $f$.)

%253+

\spheader 253Yg Misalkan $(X,\Sigma,\mu)$ dan $(Y,\Tau,\nu)$
ruang-ruang ukur, dan
$\lambda_0$ ukuran produk primitif pada $X\times Y$. Jika $f$ suatu
fungsi yang terintegralkan terhadap $\lambda_0$, tuliskan
$f_x(y)=f(x,y)$ kapan pun nilai ini terdefinisi. Tunjukkan bahwa kita
mempunyai suatu pemetaan $x\mapsto f_x^{\ssbullet}$ dari suatu subhimpunan
koterabaikan
$D_0$ dari $X$ ke $L^1(\nu)$. Tunjukkan bahwa pemetaan ini merupakan
fungsi yang terintegralkan menurut Bochner, sebagaimana didefinisikan dalam 253Yf,
%new 2008
dan bahwa integral Bochner-nya adalah $g^{\ssbullet}$, dengan
$g(y)=\int f(x,y)\mu(dx)$ kapan pun nilai ini terdefinisi; khususnya,
$\int g\,d\nu=\int f\,d\lambda_0$.
%253Yf, 253+

\spheader 253Yh Misalkan $(X,\Sigma,\mu)$ dan $(Y,\Tau,\nu)$
ruang-ruang ukur, dan
andaikan bahwa $\phi$ suatu fungsi dari $X$ ke suatu subhimpunan separabel
dari $L^1(\nu)$ yang terukur dalam arti bahwa
$\phi^{-1}[G]\in\Sigma$ untuk setiap himpunan terbuka
$G\subseteq L^1(\nu)$. Tunjukkan bahwa
terdapat suatu fungsi terukur terhadap $\Lambda$, $f$, dari $X\times Y$
ke $\Bbb
R$, dengan $\Lambda$ domain ukuran produk c.l.d.\ pada
$X\times Y$, sedemikian sehingga $\phi(x)=f_x^{\ssbullet}$ untuk setiap
$x\in X$, dengan menuliskan $f_x(y)=f(x,y)$ untuk $x\in X$, $y\in Y$.
%253+

\spheader 253Yi Misalkan $(X,\Sigma,\mu)$ dan $(Y,\Tau,\nu)$ ruang-ruang
ukur, dan $\lambda$ ukuran produk c.l.d.\ pada $X\times Y$. Tunjukkan
bahwa 253Yg memberikan identifikasi kanonik antara $L^1(\lambda)$
dan $L^1_{L^1(\nu)}(\mu)$.
%253Yg, 253Yh, 253+

\spheader 253Yj Misalkan $(X,\Sigma,\mu)$ dan $(Y,\Tau,\nu)$ ruang-ruang
ukur lengkap yang ditentukan secara lokal, dengan ukuran produk c.l.d.\ $\lambda$.
(i) Andaikan bahwa $K\in\eusm L^2(\lambda)$,
$f\in\eusm L^2(\mu)$. Tunjukkan bahwa
$h(y)=\int K(x,y)f(x)dx$ terdefinisi untuk hampir setiap
$y\in Y$ dan bahwa $h\in\eusm L^2(\nu)$. ({\it Petunjuk\/}: untuk
melihat bahwa $h$ terdefinisi hampir di mana-mana, pertimbangkan
$\int_{E\times F}K(x,y)f(x)d(x,y)$
untuk $\mu E$, $\nu F<\infty$; untuk melihat bahwa
$h\in\eusm L^2$, pertimbangkan
$\int h\times g$ dengan $g\in\eusm L^2(\nu)$.) (ii) Tunjukkan bahwa
pemetaan $f\mapsto h$ bersesuaian dengan suatu operator linear terbatas
$T_K:L^2(\mu)\to L^2(\nu)$. (iii) Tunjukkan bahwa pemetaan $K\mapsto T_K$
bersesuaian dengan suatu operator linear terbatas, bernorma paling besar
$1$, dari $L^2(\lambda)$ ke $\eurm B(L^2(\mu);L^2(\nu))$.
%253+

\spheader 253Yk Andaikan bahwa $p$, $q\in[1,\infty]$ dan bahwa
$\bover1p+\bover1q=1$, dengan menafsirkan $\bover1{\infty}$ sebagai $0$
seperti biasa. Misalkan $(X,\Sigma,\mu)$, $(Y,\Tau,\nu)$ ruang-ruang ukur
lengkap yang ditentukan secara lokal dengan ukuran produk c.l.d.\ $\lambda$.
Tunjukkan bahwa gagasan-gagasan 253Yj dapat digunakan untuk mendefinisikan
suatu operator linear terbatas, bernorma paling besar $1$, dari
$L^p(\lambda)$ ke $\eurm B(L^q(\mu);L^p(\nu))$.
%253Yj, 253+

\spheader 253Yl Dalam 253Xc, andaikan bahwa $W$ suatu kisi Banach.
Tunjukkan bahwa pernyataan-pernyataan berikut ekuivalen: (i) $Tu\ge 0$
kapan pun $u\in L^1(\lambda)$ dan $u\ge0$; (ii)
$\phi(u_1,\ldots,u_n)\ge 0$ kapan pun
$u_i\ge 0$ dalam $L^1(\mu_i)$ untuk setiap $i\le n$.
%253G, 253Xc
}%end of exercises

\endnotes{
\Notesheader{253} Di seluruh argumen utama bagian ini, saya telah
menuliskan hasil-hasilnya dalam istilah ukuran produk c.l.d.\ tersebut; tentu saja
isomorfisme yang dicatat dalam 253Yc berarti bahwa hasil-hasil itu dapat
pula dinyatakan dalam istilah ukuran produk primitif. Gagasan
keterintegralan yang lebih terbatas terhadap ukuran produk primitif
memang merupakan gagasan yang tepat bagi gagasan-gagasan 253Yg.

Teorema 253F adalah suatu 'teorema pemetaan universal';
teorema itu menyatakan bahwa setiap operator bilinear terbatas pada
$L^1(\mu)\times
L^1(\nu)$ terfaktorkan melalui
$\otimes:L^1(\mu)\times L^1(\nu)\to
L^1(\lambda)$, setidaknya jika ruang hasilnya suatu ruang Banach.
Mudah dilihat bahwa sifat ini menentukan pasangan
$(L^1(\lambda),\otimes)$ hingga isomorfisme ruang Banach, dalam arti
berikut: jika $V$ suatu ruang Banach, dan
$\psi:L^1(\mu)\times L^1(\nu)\to V$ suatu operator bilinear terbatas
sedemikian sehingga untuk setiap operator bilinear terbatas $\phi$ dari
$L^1(\mu)\times L^1(\nu)$ ke sebarang ruang Banach $W$ terdapat suatu
operator linear terbatas yang unik $T:V\to W$ sedemikian sehingga
$T\psi = \phi$ dan $\|T\|=\|\phi\|$, maka terdapat suatu isomorfisme
isometrik ruang Banach $S:L^1(\lambda)\to V$ sedemikian sehingga
$S\otimes=\psi$. Tentu saja terdapat suatu
teori umum operator bilinear antara ruang-ruang Banach; dalam bahasa
teori ini, $L^1(\lambda)$ adalah, atau isomorfik dengan,
'produk tensor projektif' dari $L^1(\mu)$ dan $L^1(\nu)$. Untuk
pengantar mengenai
subjek ini, lihat {\smc Defant \& Floret 93}, \S I.3, atau
{\smc Semadeni 71}, \S20.
Barangkali perlu saya tekankan, demi mereka yang
belum pernah menjumpai produk tensor, bahwa teorema ini
khusus bagi ruang-ruang $L^1$. Meskipun beberapa gagasan yang sama dapat
diterapkan pada ruang-ruang fungsi lain (lihat 253Xe-253Xg), tidak ada kelas
lain tempat 253F berlaku.

Terdapat pula teori produk tensor kisi Banach, tetapi saya rasa kita
belum benar-benar siap untuk itu (teori ini memerlukan gagasan-gagasan
umum tentang ruang linear terurut yang hendak saya tunda sampai Bab 35
dalam jilid berikutnya). Namun demikian,
253G menunjukkan bahwa pengurutan, dan karena itu struktur kisi Banach,
dari
$L^1(\lambda)$ ditentukan oleh pengurutan $L^1(\mu)$ dan
$L^1(\nu)$ serta pemetaan
$\otimes:L^1(\mu)\times L^1(\nu)\to L^1(\lambda)$.

Operator-operator ekspektasi bersyarat yang dideskripsikan dalam 253H
sangat penting, terutama karena dalam konteks khusus ini kita mempunyai
realisasi operator ekspektasi bersyarat sebagai suatu fungsi $P_0$
dari $\eusm L^1(\lambda)$ ke $\eusm L^1(\lambda\restr\Lambda_1)$, bukan
sekadar sebagai suatu fungsi dari
$L^1(\lambda)$ ke $L^1(\lambda\restr\Lambda_1)$, seperti dalam 242J.
Sebagaimana dideskripsikan di sini, $P_0(f+f')$ tidak harus sama,
dalam arti ketat, dengan $P_0f+P_0f'$; domainnya dapat lebih besar.
Namun, dalam penerapan,
kita mungkin bersedia membatasi perhatian pada ruang linear
$\eusm U$ dari fungsi-fungsi terbatas yang terukur terhadap
$\Sigma\tensorhat\Tau$ dan didefinisikan di mana-mana pada $X\times Y$,
sehingga $P_0$ menjadi operator dari $\eusm U$ ke dirinya sendiri
(lihat 252P).
}%end of notes

\discrpage
